% from main.tex
%%%%%%%%%%  scan idx 143  |  folio 130  %%%%%%%%%%
\textbf{Corollaire 4.19.1.} \textbf{a)} Le foncteur
\[
  \mathrm L f^*:\Dpx^-({}_A S)\longrightarrow \Dpx^-({}_A, S')
\]
induit un foncteur
\[
  \mathrm L f^*:\Dpx^-({}_A S)_{\mathrm{parf}}
  \longrightarrow
  \Dpx^-({}_A, S')_{\mathrm{parf}}.
\]

\textbf{b)} Si \(A\) est commutatif, le foncteur
\[
  \overset{\mathrm L}{\otimes}_A:
  \Dpx^-(S)\times_S\Dpx^-(S)
  \longrightarrow \Dpx^-(S)
\]
induit un foncteur
\[
  \overset{\mathrm L}{\otimes}_A:
  \Dpx^-(S)_{\mathrm{parf}}\times_S\Dpx^-(S)_{\mathrm{parf}}
  \longrightarrow \Dpx^-(S)_{\mathrm{parf}}.
\]

\textbf{Remarque 4.19.2.} On peut préciser (4.19) en termes d'amplitude parfaite :
\[
  \perfamp(E)\subset[a,b]\quad\text{et}\quad
  \perfamp(E')\subset[a',b']
  \quad\Rightarrow\quad
  \perfamp(E'\overset{\mathrm L}{\otimes}_A E)\subset[a+a',b+b'].
\]
En particulier :
\[
  \perfamp(E)\subset[a,b]
  \quad\Rightarrow\quad
  \perfamp(\mathrm L f^*(E))\subset[a,b].
\]

\section*{5. Tor-dimension finie et perfection}

\textbf{5.0.} Le lecteur se sera sûrement posé la question : que manque-t-il à un complexe pseudo-cohérent \(E\) pour être parfait ? Suffit-il, par exemple, que \(E\) soit à cohomologie bornée ? La réponse est non. Prenons en effet pour \(\mathcal C\) la catégorie des modules sur un anneau commutatif noethérien \(A\), et pour \(\mathcal C_0\) la sous-catégorie pleine des \(A\)-modules projectifs de type fini. N'importe quel \(A\)-module de type fini \(E\) est pseudo-cohérent et à cohomologie bornée, mais \(E\) n'est parfait que si \(E\) est en outre de dimension projective finie (ou, ce qui revient au même, de tor-dimension finie). Cet exemple suggère que la différence entre un complexe pseudo-cohérent et un complexe parfait réside dans une question de dimension cohomologique. C'est ce que nous allons examiner maintenant
